\begin{abstract}
This dated working manuscript is an incomplete, conditional
human-readable reconstruction of the classical Riemann-hypothesis
specialization of the Alp\"oge--Claude formal argument. It is not an
unconditional proof of the Riemann hypothesis and has not been
independently verified as a mathematical manuscript. We show that the
Riemann hypothesis follows if, for every $\delta>0$, the error $G(x)$
in a smoothed prime sum satisfies
$G(x)\le C_\delta x^{1/2+\delta}$ on $[Y,Y^{3/2}]$ for arbitrarily
large $Y$. In Levent Alp\"oge's adaptation of Stepanov's method for
curves over finite fields, divisibility at inert primes along a chain
of theta lifts plays the role of high-order vanishing at rational
points. For a sum of length $x$, Claude proposed taking about
$\sqrt{x}$ lifts \citep[p.~1]{alpoge-remark}. We compare this
divisibility with the heights of the period maps using one set of
weights and one shift for all lifts at each mesh point. The required
uniform arithmetic estimate and several named inputs remain open.
Research on their reconstruction is ongoing.
\end{abstract}
